\documentclass{article}

\usepackage{fancyhdr}
\usepackage{extramarks}
\usepackage{amsmath}
\usepackage{amsthm}
\usepackage{amsfonts}
\usepackage{tikz}
\usepackage[plain]{algorithm}
\usepackage{algpseudocode}

\usetikzlibrary{automata,positioning}

%
% Basic Document Settings
%

\topmargin=-0.45in
\evensidemargin=0in
\oddsidemargin=0in
\textwidth=6.5in
\textheight=9.0in
\headsep=0.25in

\linespread{1.1}

\pagestyle{fancy}
\lhead{\hmwkAuthorName}
\rhead{\hmwkClass:\ \hmwkTitle}
\cfoot{\thepage}

\renewcommand\headrulewidth{0.4pt}
\renewcommand\footrulewidth{0.4pt}

\setlength\parindent{0pt}

%
% Create Problem Sections
%

% Section variable
\newcommand{\currentSection}{3}

\newcommand{\enterProblemHeader}[1]{
    \nobreak\extramarks{}{Problem #1 continued on next page\ldots}\nobreak{}
    \nobreak\extramarks{Problem #1 (continued)}{Problem #1 continued on next page\ldots}\nobreak{}
}

\newcommand{\exitProblemHeader}[1]{
    \nobreak\extramarks{Problem #1 (continued)}{Problem #1 continued on next page\ldots}\nobreak{}
    \nobreak\extramarks{Problem #1}{}\nobreak{}
}

\newenvironment{homeworkProblem}[1]{
    \def\tempProblemNum{\currentSection.#1}
    \section*{Problem \tempProblemNum}
    \setcounter{partCounter}{1}
    \enterProblemHeader{\tempProblemNum}
}{
    \exitProblemHeader{\tempProblemNum}
}

\setcounter{secnumdepth}{0}
\newcounter{partCounter}
\newcounter{homeworkProblemCounter}
\setcounter{homeworkProblemCounter}{1}
\nobreak\extramarks{Problem \arabic{homeworkProblemCounter}}{}\nobreak{}

%
% Homework Problem Environment
%
% This environment takes an optional argument. When given, it will adjust the
% problem counter. This is useful for when the problems given for your
% assignment aren't sequential. See the last 3 problems of this template for an
% example.
%
\renewenvironment{homeworkProblem}[1]{
    \def\tempProblemNum{#1} 
    \section*{Problem \currentSection.#1}
    \setcounter{partCounter}{1}
    \enterProblemHeader{\currentSection.#1}
}{
    \exitProblemHeader{\currentSection.\tempProblemNum}
}

%
% Homework Details
%   - Title
%   - Due date
%   - Class
%   - Section/Time
%   - Instructor
%   - Author
%

\newcommand{\hmwkTitle}{Problem Set 6}
\newcommand{\hmwkDueDate}{10 March 2026}
\newcommand{\hmwkClass}{Foundations of Advanced Mathematics}
\newcommand{\hmwkClassTime}{MATH355}
\newcommand{\hmwkClassInstructor}{Dr. Anthony Bosman}
\newcommand{\hmwkAuthorName}{\textbf{Micah Leiterman}}

%
% Title Page
%

\title{
    \vspace{2in}
    \textmd{\textbf{\hmwkClass:\\\ \hmwkTitle}}\\
    \normalsize\vspace{0.1in}\small{Due\ on\ \hmwkDueDate}\\
    \vspace{0.1in}\large{\textit{\hmwkClassInstructor\ \hmwkClassTime}}
    \vspace{3in}
}

\author{\hmwkAuthorName}
\date{}

\renewcommand{\part}[1]{\textbf{\large Part \Alph{partCounter}}\stepcounter{partCounter}\\}

%
% Various Helper Commands
%

% Useful for algorithms
\newcommand{\alg}[1]{\textsc{\bfseries \footnotesize #1}}

% For derivatives
\newcommand{\deriv}[1]{\frac{\mathrm{d}}{\mathrm{d}x} (#1)}

% For partial derivatives
\newcommand{\pderiv}[2]{\frac{\partial}{\partial #1} (#2)}

% Integral dx
\newcommand{\dx}{\mathrm{d}x}

% Alias for the Solution section header
\newcommand{\solution}{\textbf{\large Solution}}

% Probability commands: Expectation, Variance, Covariance, Bias
\newcommand{\E}{\mathrm{E}}
\newcommand{\Var}{\mathrm{Var}}
\newcommand{\Cov}{\mathrm{Cov}}
\newcommand{\Bias}{\mathrm{Bias}}

\begin{document}

\maketitle

\pagebreak

\begin{homeworkProblem}{17a}
    \begin{center}
    \begin{tabular}{l | c c c c c c c}
        \(\mathbf{+}\) & \(\mathbf{0}\) & \(\mathbf{1}\) & \(\mathbf{2}\) & \(\mathbf{3}\) & \(\mathbf{4}\) & \(\mathbf{5}\) & \(\mathbf{6}\) \\ \hline
        \(\mathbf{0}\) & 0 & 1 & 2 & 3 & 4 & 5 & 6 \\ 
        \(\mathbf{1}\) & 1 & 2 & 3 & 4 & 5 & 6 & 0 \\ 
        \(\mathbf{2}\) & 2 & 3 & 4 & 5 & 6 & 0 & 1 \\ 
        \(\mathbf{3}\) & 3 & 4 & 5 & 6 & 0 & 1 & 2 \\ 
        \(\mathbf{4}\) & 4 & 5 & 6 & 0 & 1 & 2 & 3 \\ 
        \(\mathbf{5}\) & 5 & 6 & 0 & 1 & 2 & 3 & 4 \\ 
        \(\mathbf{6}\) & 6 & 0 & 1 & 2 & 3 & 4 & 5 \\ 
    \end{tabular}
    \end{center}
\end{homeworkProblem}

\begin{homeworkProblem}{17a}
    \begin{center}
    \begin{tabular}{l | c c c c c c c}
        \(\mathbf{\times}\) & \(\mathbf{0}\) & \(\mathbf{1}\) & \(\mathbf{2}\) & \(\mathbf{3}\) & \(\mathbf{4}\) & \(\mathbf{5}\) & \(\mathbf{6}\) \\ \hline
        \(\mathbf{0}\) & 0 & 0 & 0 & 0 & 0 & 0 & 0 \\ 
        \(\mathbf{1}\) & 0 & 1 & 2 & 3 & 4 & 5 & 6 \\ 
        \(\mathbf{2}\) & 0 & 2 & 4 & 6 & 1 & 3 & 5 \\ 
        \(\mathbf{3}\) & 0 & 3 & 6 & 2 & 5 & 1 & 4 \\ 
        \(\mathbf{4}\) & 0 & 4 & 1 & 5 & 2 & 6 & 3 \\ 
        \(\mathbf{5}\) & 0 & 5 & 3 & 1 & 6 & 4 & 2 \\ 
        \(\mathbf{6}\) & 0 & 6 & 5 & 4 & 3 & 2 & 1 \\ 
    \end{tabular}
    \end{center}
\end{homeworkProblem}

\begin{homeworkProblem}{17b}
   For $[a \in \{1,2,...,6\}$, a has an additive inverse for all values in the set: \\
   \[ \{0, 1, 2, 3, 4, 5, 6\} \]
   \[ \{0, 1, 2, 3, 4, 5, 6\} \]
    
\end{homeworkProblem}

\begin{homeworkProblem}{17c}
   For $[a \in \{1,2,...,6\}$, a has a multiplicative inverse for all values in the set because 7 is prime: \\
   \[ \{0, 1, 2, 3, 4, 5, 6\} \]
   \[ \{0, 1, 2, 3, 4, 5, 6\} \]

\end{homeworkProblem}

\begin{homeworkProblem}{17d}
    \begin{center}
    \begin{tabular}{l | c c c c c c c c}
        \(\mathbf{\times}\) & \(\mathbf{0}\) & \(\mathbf{1}\) & \(\mathbf{2}\) & \(\mathbf{3}\) & \(\mathbf{4}\) & \(\mathbf{5}\) & \(\mathbf{6}\) & \(\mathbf{7}\) \\ \hline
        \(\mathbf{0}\) & 0 & 0 & 0 & 0 & 0 & 0 & 0 & 0 \\ 
        \(\mathbf{1}\) & 0 & 1 & 2 & 3 & 4 & 5 & 6 & 7 \\ 
        \(\mathbf{2}\) & 0 & 2 & 4 & 6 & 0 & 2 & 4 & 6 \\ 
        \(\mathbf{3}\) & 0 & 3 & 6 & 1 & 4 & 7 & 2 & 5 \\ 
        \(\mathbf{4}\) & 0 & 4 & 0 & 4 & 0 & 4 & 0 & 4 \\ 
        \(\mathbf{5}\) & 0 & 5 & 2 & 7 & 4 & 1 & 6 & 3 \\ 
        \(\mathbf{6}\) & 0 & 6 & 4 & 2 & 0 & 6 & 4 & 2 \\ 
        \(\mathbf{7}\) & 0 & 7 & 6 & 5 & 4 & 3 & 2 & 1 \\ 
    \end{tabular}
    \end{center}

    
    For \(a \in \{0, 1, 2, 3, 4, 5, 6, 7\}\), \(a\) has an additive inverse for all values in the set:
    \[ \{0, 1, 2, 3, 4, 5, 6, 7\} \]

  
    For \(a \in \{0, 1, 2, 3, 4, 5, 6, 7\}\), \(a\) has a multiplicative inverse only when \(\gcd(a, 8) = 1\). These values are:
    \[ \{1, 3, 5, 7\} \]
\end{homeworkProblem}

\pagebreak

\begin{homeworkProblem}{19}
   We are looking for $1!+2!+3!+...+100!$ (mod 12). \\
   We can disregard all factorials 12 and above because they are 12 times something so they are $\equiv$ 0 (mod 12).
   All the factorials 4 and above contain $3\times4$ which means they are $\equiv$ 0 (mod 12).\\
   $1+2+6 \equiv 9$ (mod 12).
   
   
\end{homeworkProblem}

\pagebreak

\begin{homeworkProblem}{22a}
    To find $13^8$(mod 7), we first look at the base 13. \\
    $13 = 7+6$, which means $13 \equiv -1$ (mod 7). \\
    Substituting in we get $(-1)^8$ which is $\equiv 1$ (mod 7).
\end{homeworkProblem}

\begin{homeworkProblem}{22b}
    To find $2^{2021}$(mod 7), we first look at the pattern for powers of 2. \\
    \[2^1 \equiv 2, 2^2 \equiv 4, 2^3 \equiv 1.\]
    We can factor the exponent 2021 into \[(3\times336) +5\]
    giving us $2^{(3\times336) +5}$.
    \[2^{(3\times336 +5)}\]
    \[\equiv 2^{3^{336}} \times 2^{3+2}\]
    \[\equiv 1 \times 2^2\]
    \[\equiv 4\]
    This means that the last digit is 4.
\end{homeworkProblem}

\begin{homeworkProblem}{22c}
    The last digit of $7^{2021}$ is the same thing as saying $7^{2021}$ (mod 10). \\

    $7^2 = 49 \equiv -1$ (mod 10).\\
    We rewrite $7^{2021}$ as $7^{2^{1010} + 1}$.
    \[ 7^{2^{1010} + 1}\]
    \[ =(-1)^{1010} \times 7^1\]
    \[\equiv 7\]
\end{homeworkProblem}

\begin{homeworkProblem}{22d}
    The last two digits of $11^{100}$ is the same thing as saying $11^{100}$ (mod 100). \\

    \[11^1 \equiv 11, 11^2 \equiv 21, 11^3 \equiv 31,... 11^{10}\equiv 01, 11^{11} \equiv 11\].
    This means that 
    \[11^{100}\]
    \[=11^{10^{10}}\]
    \[\equiv01\].
   
\end{homeworkProblem}
\pagebreak

\begin{homeworkProblem}{23a}
    Find an integer \(x\) such that:
    \[
        \begin{cases}
            x \equiv 2 \pmod 3 \\
            x \equiv 3 \pmod 4
        \end{cases}
    \]
    From \(x \equiv 2 \pmod 3\), we have \(x = 3k + 2\) for some integer \(k\). 
    Substituting this into the second congruence:
    \[
        \begin{split}
            3k + 2 &\equiv 3 \pmod 4 \\
            3k &\equiv 1 \pmod 4
        \end{split}
    \]
    The modular inverse of \(3 \pmod 4\) is 3. Multiplying both sides by 3:
    \[
        k \equiv 3 \pmod 4
    \]
    Substituting \(k = 3\) back into the expression for \(x\):
    \[
        x = 3(3) + 2 = 11
    \]
    The number 11 satisfies both congruences.
\end{homeworkProblem}

\begin{homeworkProblem}{23b}
    Find an integer x such that:
    \[
        \begin{cases}
            x \equiv 2 \pmod 3 \\
            x \equiv 4 \pmod 5 \\
            x \equiv 1 \pmod 8
        \end{cases}
    \]
    From the first congruence, \(x = 3k + 2\). Substituting into the second:
    \[
        \begin{split}
            3k + 2 &\equiv 4 \pmod 5 \\
            3k &\equiv 2 \pmod 5 \\
            k &\equiv 4 \pmod 5 \implies k = 5j + 4
        \end{split}
    \]
    Substituting back into our expression for \(x\):
    \[
        x = 3(5j + 4) + 2 = 15j + 14
    \]
    Now, substituting this into the third congruence:
    \[
        \begin{split}
            15j + 14 &\equiv 1 \pmod 8 \\
            -j + 6 &\equiv 1 \pmod 8 \\
            -j &\equiv -5 \pmod 8 \\
            j &\equiv 5 \pmod 8
        \end{split}
    \]
    Plugging \(j = 5\) into the expression for \(x\):
    \[
        x = 15(5) + 14 = 89
    \]
    The number 89 satisfies all three conditions.
\end{homeworkProblem}
\pagebreak

\begin{homeworkProblem}{23c}
    This number cannot exist. We will prove its nonexistence by contradiction. That is, suppose some number $x$ is congruent to 1 (mod 3), 2 (mod 3), and 3 (mod 10). We can write $x$ as $x =3k+1$ for it to be 1 (mod 3). For it to also be 2 (mod 3), x must also be able to be represented by $x=3k +2$.
    \[\begin{split}
        3k+1 = x = 3k+2\\
        3k+1 = 3k+2\\
        1 = 2
    \end{split}\]
    We have found a contradiction, so this number cannot exist.
\end{homeworkProblem}
\begin{homeworkProblem}{23d}
    Find an integer \(x\) such that:
    \[
        \begin{cases}
            x \equiv 1 \pmod 2 \\
            x \equiv 2 \pmod 3 \\
            x \equiv 3 \pmod 5 \\
            x \equiv 4 \pmod{11}
        \end{cases}
    \]
    We combine the congruences iteratively. 
    From \(x \equiv 1 \pmod 2\) and \(x \equiv 2 \pmod 3\), we find \(x = 6j + 5\).
    Substituting into the third congruence:
    \[
        \begin{split}
            6j + 5 &\equiv 3 \pmod 5 \\
            j &\equiv 3 \pmod 5 \implies j = 5m + 3
        \end{split}
    \]
    Thus, \(x = 6(5m + 3) + 5 = 30m + 23\). Substituting into the final congruence:
    \[
        \begin{split}
            30m + 23 &\equiv 4 \pmod{11} \\
            8m + 1 &\equiv 4 \pmod{11} \\
            8m &\equiv 3 \pmod{11}
        \end{split}
    \]
    Multiplying by the modular inverse of \(8 \pmod{11}\), which is \(7\):
    \[
        m \equiv 21 \equiv 10 \pmod{11}
    \]
    Plugging \(m = 10\) back into the expression for \(x\):
    \[
        x = 30(10) + 23 = 323
    \]
    The number 323 satisfies all four conditions.
\end{homeworkProblem}
\pagebreak

\begin{homeworkProblem}{24a}
    To find $5^{50}$ (mod 7), we use Fermat's Little Theorem that states $a^{p-1} \equiv 1$ for a prime $p$ and any integer $a$. We can apply it here because 7 is prime.\\
    We start with factoring $50$ into $(6\times8) + 2$.\\
    From there:
    \[
        \begin{split}
        &=5^{50}\\
        &=5^{(6\times8)+2}\\
        &=5^{(6^8)+2}\\
        &=5^{(6^8)}\times5^2\\
        &=({5^{6}})^8 \times 5^2 \\
        &=1^8 \times 5^2 \text{ by FLT $5^6 \equiv 1\pmod{7}$}\\
        &=5^2\\
        &=25 \equiv 4 \pmod{7}.
        \end{split}
    \]
    Therefore, $5^{50} \equiv 4 \pmod{7}$.
    
\end{homeworkProblem}

\begin{homeworkProblem}{24b}
    We wish to calculate \(12^{101} \pmod 5\) using Fermat's Little Theorem.
    
    First, we simplify the base mod 5:
    \[
        12 \equiv 2 \pmod 5
    \]
    Since 5 is prime and relatively prime with 2, Fermat's Little Theorem states:
    \[
        2^{5-1} = 2^4 \equiv 1 \pmod 5
    \]
    Next, we factor the exponent:
    \[
        101 = 25(4) + 1 \implies 101 \equiv 1 \pmod 4
    \]
    Applying these reductions:
    \[
        \begin{split}
            &\equiv 2^{101} \\
            &\equiv (2^4)^{25} \times 2^1 \pmod{5}\\
            &\equiv (1)^{25} \times 2 \\
            &\equiv 2 \pmod 5
        \end{split}
    \]
    Therefore, the answer is $2$.
\end{homeworkProblem}
\begin{homeworkProblem}{24c}
    We wish to calculate \(3^{2026} \pmod 5\) using Fermat's Little Theorem.
    
    Since 5 is a prime number and \(\gcd(3, 5) = 1\), Fermat's Little Theorem states:
    \[
        3^{5-1} = 3^4 \equiv 1 \pmod 5
    \]
    Next, we divide the exponent 2026 by 4 to find the remainder:
    \[
        2026 = 506(4) + 2 \implies 2026 \equiv 2 \pmod 4
    \]
    Substituting this into the power expression:
    \[
        \begin{split}
            3^{2026} &\equiv 3^2 \pmod 5 \\
            3^{2026} &\equiv 9 \pmod 5 \\
            3^{2026} &\equiv 4 \pmod 5
        \end{split}
    \]
    The final answer is 4.
\end{homeworkProblem}

\begin{homeworkProblem}{24d}
    We wish to calculate \(3^{2026} \pmod 7\) using Fermat's Little Theorem.
    
    Since 7 is a prime number and \(\gcd(3, 7) = 1\), Fermat's Little Theorem states:
    \[
        3^{7-1} = 3^6 \equiv 1 \pmod 7
    \]
    Next, we divide the exponent 2026 by 6 to find the remainder:
    \[
        2026 = 337(6) + 4 \implies 2026 \equiv 4 \pmod 6
    \]
    Substituting this into the power expression:
    \[
        \begin{split}
            3^{2026} &\equiv 3^4 \pmod 7 \\
            3^{2026} &\equiv 81 \pmod 7 \\
            3^{2026} &\equiv 4 \pmod 7
        \end{split}
    \]
    The final answer is 4.
\end{homeworkProblem}
\begin{homeworkProblem}{24e}
    We wish to calculate \(3^{2026} \pmod{11}\) using Fermat's Little Theorem.
    
    Since 11 is a prime number and \(\gcd(3, 11) = 1\), Fermat's Little Theorem states:
    \[
        3^{11-1} = 3^{10} \equiv 1 \pmod{11}
    \]
    Next, we divide the exponent 2026 by 10 to find the remainder:
    \[
        2026 = 202(10) + 6 \implies 2026 \equiv 6 \pmod{10}
    \]
    Substituting this into the power expression:
    \[
        \begin{split}
            3^{2026} &\equiv 3^6 \pmod{11} \\
            3^3 &= 27 \equiv 5 \pmod{11} \\
            3^6 &= (3^3)^2 \equiv 5^2 \equiv 25 \pmod{11} \\
            25 &\equiv 3 \pmod{11}
        \end{split}
    \]
    The final answer is 3.
\end{homeworkProblem}
\pagebreak

    \begin{homeworkProblem}{26}
    We wish to prove that \(42 \mid (n^7 - n)\) for all \(n \in \mathbb{N}\). 
    Note that \(42 = 2 \cdot 3 \cdot 7\).
    
    By Fermat's Little Theorem, \(n^p \equiv n \pmod p\) for any prime \(p\). For \(p = 7\):
    \[ n^7 \equiv n \pmod 7 \implies n^7 - n \equiv 0 \pmod 7 \]
    By Fermat's Little Theorem, \(n^3 \equiv n \pmod 3\). We can write:
    \[ n^7 - n = n(n^6 - 1) = n(n^2 - 1)(n^4 + n^2 + 1) = (n-1)n(n+1)(n^4 + n^2 + 1) \]
    Since \((n-1)n(n+1)\) is a product of three consecutive integers, it is divisible by \(3! = 6\). Thus, \(n^7 - n \equiv 0 \pmod 3\).
    
    Since \(n^7 - n = n(n-1)(\dots)\), and the product of two consecutive integers is always even, \(n^7 - n \equiv 0 \pmod 2\).
    
    Because \(n^7 - n\) is divisible by 2, 3, and 7, and \(\gcd(2,3,7) = 1\), it is divisible by their product \(2 \cdot 3 \cdot 7 = 42\).
    \[ 42 \mid (n^7 - n) \]
\end{homeworkProblem}
\pagebreak
\begin{homeworkProblem}{29}
\textbf{a. Verify $d=5$ as a private key}
    Given public key \((n, e) = (85, 13)\):

    Factors of \(n=85\) are \(p=5, q=17\). 
     \(\phi = (5-1)(17-16) = 64\).
    The private key is $d$ such that \(ed \equiv 1 \pmod{\phi}\):
    \[ 13 \times 5 = 65 \equiv 1 \pmod{64} \]

    \textbf{b. Encryption of \(m=7\)}
    \[ c = m^e \pmod n = 7^{13} \pmod{85} \]
    Look at the pattern of powers of 7. 
    \[ 7^1 \equiv 7, \quad 7^2 \equiv 49, \quad 7^4 \equiv 21, \quad 7^8 \equiv 16 \]
    \[ c \equiv 16 \cdot 21 \cdot 7 = 2352 \equiv 57 \pmod{85} \]

    \textbf{c. Decryption of \(c=57\)}
    \[ m = c^d \pmod n = 57^5 \pmod{85} \]
    \[ 57^2 \equiv 19, \quad 57^4 \equiv 21 \]
    \[ m \equiv 21 \cdot 57 = 1197 \equiv 7 \pmod{85} \]
    The decrypted message is 7.
\end{homeworkProblem}
\begin{homeworkProblem}{21}
    Let \( n = \sum_{i=0}^{k} d_i 10^i \). We wish to show that \( 11 \mid n \iff 11 \mid \sum_{i=0}^{k} (-1)^i d_i \).

    Recall that \( 10 \equiv -1 \pmod{11} \). By the properties of modular exponentiation:
    \[ 10^a \equiv (-1)^a \pmod{11} \]
    
    Applying this to the decimal expansion of \( n \):
    \[
        \begin{split}
            n &= d_k 10^k + d_{k-1} 10^{k-1} + \dots + d_1 10^1 + d_0 10^0 \\
            n &\equiv d_k (-1)^k + d_{k-1} (-1)^{k-1} + \dots + d_1 (-1)^1 + d_0 (1) \pmod{11} \\
            n &\equiv \sum_{i=0}^{k} (-1)^i d_i \pmod{11}
        \end{split}
    \]
    
    An integer is divisible by 11 iff it is congruent to \( 0 \pmod{11} \). 
    Since \( n \) and its alternating digit sum are congruent modulo 11, one is divisible by 11 if and only if the other is.
\end{homeworkProblem}
\end{document}
